\setcounter{section}{0}
\renewcommand{\thesection}{\Roman{section}}

\begin{center}
{\LARGE\bfseries Impact of Network Topology on Distributed Inference:\\
Circulant Graphs as a Replacement for 3D Torus\par}
\vspace{1.0em}
\begin{tabular}{cc}
Nikhil Vemuri & Brian Zhang \\
\texttt{nvemuri@mit.edu} & \texttt{zhangbr@mit.edu}
\end{tabular}
\end{center}
\vspace{1.0em}

\section{Key Insights and Innovations}
This project couples AccelForge mapping with a topology-aware replay model for
64-chip accelerator inference. The central design insight is that
$C(64,\{1,5,17\})$ preserves the six-link-per-chip budget of a TPU-v4-like 3D
torus while decomposing into three edge-disjoint Hamiltonian cycles, which
spreads full-chip reductions over independent rings. This lowers the predicted
AllReduce and ReduceScatter bottleneck load by about $2.29\times$ relative to
the torus dimension-wise schedule.

Across the evaluated small, tall, square, and MoE expert-FFN workloads, the
circulant graph wins every case. The results show that serialization
bottlenecks, not first-byte latency, dominate these communication-heavy
mappings. The feedback loop also converges in two iterations in the evaluated
runs, suggesting that topology mainly changes the cost of a stable collective
plan for these workloads.

\section{Contributions}
Nikhil Vemuri and Brian Zhang jointly designed the topology study, developed
the AccelForge-to-network replay workflow, ran the matmul and MoE experiments,
analyzed the resulting latency and energy trends, and wrote the final report.
Nikhil focused on topology modeling, collective routing, and MoE and full architecture workload
experiments. Brian focused on AccelForge integration, congestion and latency modeling, and matmul workload experiments.

\section{Artifact}
The project artifact is the repository
\texttt{https://github.com/WubbLord/network-topology}. The main report is
\texttt{report/main.tex}; the technical summary is this file,
\texttt{report/technical\_summary.tex}.
